% Chapter 1: Introduction
% Target: 12-15 pages (3,000-4,000 words)

\section{Research Background}

The intersection of artificial intelligence and finance has undergone transformative changes over the past decade. Traditional quantitative trading strategies, once the domain of statistical arbitrage and technical analysis, have evolved to incorporate increasingly sophisticated machine learning techniques. This evolution, however, has reached an inflection point with the emergence of large language models (LLMs) capable of sophisticated reasoning about financial markets.

\subsection{The Rise of Machine Learning in Finance}

Machine learning has become ubiquitous in quantitative finance, with applications ranging from price prediction to portfolio optimization. The appeal is straightforward: financial markets generate vast amounts of data, and ML models excel at identifying patterns within such datasets. Landmark achievements include:

\begin{itemize}
    \item LSTM-based models for sequence prediction \cite{hochreiter1997lstm}
    \item Gradient boosting methods for feature-rich environments \cite{ke2017lightgbm}
    \item Reinforcement learning for optimal execution
\end{itemize}

Despite these successes, ML approaches in finance face fundamental challenges:

\textbf{Data Limitations}: ML models require large amounts of labeled training data. In finance, labels (future returns) are inherently noisy, and the signal-to-noise ratio is extremely low.

\textbf{Non-stationarity}: Financial time series exhibit regime changes---the statistical properties that hold in one period may not hold in another. This violates the i.i.d. assumption underlying most ML algorithms.

\textbf{Interpretability Gap}: Complex ML models are often ``black boxes,'' making it difficult for practitioners to understand and trust their predictions.

\subsection{The LLM Paradigm Shift}

The release of GPT-4 and subsequent large language models has opened new possibilities for financial applications. Unlike traditional ML models that rely solely on numerical features, LLMs can:

\begin{enumerate}
    \item Process unstructured text data (news, earnings calls, social media)
    \item Engage in chain-of-thought reasoning about market dynamics
    \item Generate human-interpretable explanations for predictions
\end{enumerate}

Recent work has demonstrated the potential of LLMs in financial tasks, including sentiment analysis \cite{kirtac2024sentiment}, financial question answering, and trading signal generation \cite{yang2023fingpt}.

\section{Problem Statement}

While both ML and LLM approaches have shown promise in quantitative trading, there lacks a systematic framework for comparing their strengths and weaknesses. Key questions remain unanswered:

\begin{enumerate}
    \item \textbf{Performance}: Do LLMs achieve superior risk-adjusted returns compared to traditional ML?
    \item \textbf{Risk Control}: How do LLMs perform during high volatility and black swan events?
    \item \textbf{Cost-Efficiency}: Is the additional computational cost of LLMs justified?
\end{enumerate}

This thesis addresses these questions through the design and implementation of a Dual-Track quantitative trading framework.

\section{Research Hypotheses}

We formulate three primary hypotheses:

\begin{description}
    \item[H1] LLM-based trading strategies achieve superior risk-adjusted returns (Sharpe ratio) compared to traditional ML approaches.
    \item[H2] LLM strategies demonstrate better risk control during high volatility periods and black swan events.
    \item[H3] The computational overhead of LLMs is justified by improved risk-adjusted returns.
\end{description}

\section{Research Contributions}

This thesis makes the following contributions:

\begin{enumerate}
    \item \textbf{Methodological}: We propose the first systematic framework for fair comparison between ML and LLM approaches in quantitative trading, including standardized data pipelines, evaluation metrics, and statistical tests.

    \item \textbf{Empirical}: We provide comprehensive backtest results on two markets (CSI300 and QQQ) with six different models, generating novel insights about the trade-offs between fitting and reasoning paradigms.

    \item \textbf{Engineering}: We release an open-source implementation with full reproducibility, including trained models, backtest results, and evaluation scripts.
\end{enumerate}

\section{Thesis Structure}

The remainder of this thesis is organized as follows:

\begin{itemize}
    \item \textbf{Chapter 2}: Reviews related work in ML for trading, financial NLP, and LLM applications.
    \item \textbf{Chapter 3}: Describes the Dual-Track framework architecture and implementation.
    \item \textbf{Chapter 4}: Details the experimental setup, datasets, and evaluation methodology.
    \item \textbf{Chapter 5}: Presents quantitative results and analysis across multiple dimensions.
    \item \textbf{Chapter 6}: Discusses findings, limitations, and implications.
    \item \textbf{Chapter 7}: Concludes with summary and future directions.
\end{itemize}